\chapter{SUSTAINABLE DEVELOPMENT GOALS (SDG) AND SOCIETAL IMPACT}
\label{app:sdg_impact}

\section{Alignment with United Nations Sustainable Development Goals}
The United Nations 2030 Agenda for Sustainable Development establishes seventeen interconnected goals to address global challenges regarding education, sustainable industrial growth, and responsible consumption. The proposed multi-task speech distillation framework aligns directly with three core Sustainable Development Goals, as delineated in Table \ref{tab:sdg_alignment}.

\begin{table}[htbp]
\centering
\caption{System Alignment with UN Sustainable Development Goals}
\label{tab:sdg_alignment}
\begin{tabular}{lp{4.2cm}p{7.5cm}}
\toprule
\textbf{Goal Number} & \textbf{Designation} & \textbf{Direct System Contribution and Relevance} \\
\midrule
\textbf{SDG 4} & Quality Education & Provides highly efficient, low-latency speech transcription and tone-aware tutoring interfaces for educational institutions and learners in remote regions with limited cloud connectivity and hardware resources. \\
\addlinespace
\textbf{SDG 9} & Industry, Innovation, and Infrastructure & Advances democratized edge intelligence by replacing expensive server-side foundation models with ultra-compact, high-performance acoustic networks deployable on low-cost hardware. \\
\addlinespace
\textbf{SDG 12} & Responsible Consumption and Production & Drastically diminishes the electrical power consumption and carbon footprint of artificial intelligence inference workloads by compressing memory requirements by over 98\% (from multi-gigabyte LLMs down to 50.2 MB). \\
\bottomrule
\end{tabular}
\end{table}

\section{Societal and Accessibility Impact}
The societal value of the developed system manifests across several vital accessibility and communicative domains:
\begin{enumerate}
    \item \textbf{Assistive Interfaces for Speech Impairments:} Individuals with dysarthria, vocal tremors, or glottal motor difficulties frequently exhibit acoustic pitch shifts and erratic speech cadence that cause conventional speech engines to fail. The regularizing influence of entropic Optimal Transport acts as a semantic stabilizer, deciphering phonological content despite severe acoustic distortion.
    \item \textbf{Empathetic Mental Health Monitoring:} By decoding continuous affective states from raw vocal prosody without relying on cloud servers, the framework enables continuous, privacy-preserving monitoring of mood trajectories (such as emerging depressive or agitated episodes) in psychiatric telecare.
    \item \textbf{Bridging the Digital Divide:} Enabling offline speech recognition on commodity multi-core processors brings state-of-the-art voice computing to developing communities lacking stable high-speed internet infrastructure.
\end{enumerate}

\section{Real-Time Application Domains}
\begin{itemize}
    \item \textbf{Smart Automotive In-Cabin Systems:} Simultaneously recognizing driver navigation commands while monitoring stress or drowsiness from vocal inflection to adjust vehicle safety systems.
    \item \textbf{Automated Telephonic Emergency Triage:} Transcribing emergency 911/108 dispatch calls in real-time while estimating caller panic or trauma levels to prioritize ambulance dispatch queues.
    \item \textbf{Hands-Free Wearable Industrial Interfaces:} Providing heads-up speech recognition for manufacturing and surgical personnel in loud, acoustically challenging environments.
\end{itemize}
